% =====================================================================
%  BÖLÜM 9: Koşullu Beklenti ve Martingaller  [OPSİYONEL / İLERİ]
%  Kaynak: Williams; Durrett Böl.4; Klenke Böl.9-12; Shiryaev Böl.VII
% =====================================================================
\chapter{Koşullu Beklenti ve Martingaller}
\label{chp:martingaller}

\bolumalinti{%
  Martingal, dürüst bir oyunun matematiksel modelidir.%
}{Joseph Leo Doob, \textit{Stochastic Processes}}

% ---- Kazanımlar (TYMM) ----
\begin{kazanimlar}
Bu bölümün sonunda öğrenci:
\begin{itemize}[leftmargin=1.8em]
  \item Koşullu beklentinin Radon-Nikodym tanımını ve temel özelliklerini açıklayabilir.
  \item Martingal, süpmartingal ve altmartingal kavramlarını tanımlayabilir; örnekler verebilir.
  \item Durma zamanı ve isteğe bağlı durma teoremini ifade edebilir.
  \item Doob'un ayrışım ve yakınsama teoremlerini kullanabilir.
  \item Martingal yakınsamasının GBSK ile bağlantısını açıklayabilir.
\end{itemize}
\end{kazanimlar}

\begin{motivasyon}
Bölüm~\ref{sec:kosullu_beklenti}'de koşullu beklentiyi $L^2$ projeksiyon olarak tanımladık.
Bu bölümde \emph{bilgi akışı} boyutunu ekliyoruz: zamanla biriken gözlemler $\mathcal{F}_0 \subseteq \mathcal{F}_1 \subseteq \cdots$
bir filtrasyon oluşturur. Martingal, "dürüst oyun" sezgisini bu dil ile yakalar:
geçmişi bilince gelecekteki beklenti bugünkü değere eşit.

\medskip
\textbf{Köprü.} Bölüm~\ref{sec:kosullu_beklenti}'deki koşullu beklenti ve
Bölüm~\ref{chp:buyuk_sayilar}'daki Borel-Cantelli bu bölümün araçlarıdır.
Martingal yakınsaması, Güçlü BSK'nın zarif bir genellemesidir.
\end{motivasyon}

% =====================================================================
\section{Filtrasyon ve Uyum}
\label{sec:filtrasyon}
% =====================================================================

\begin{tanim}[Filtrasyon]
\label{def:filtrasyon}
$(\Omega,\mathcal{F},P)$ olasılık uzayı üzerinde \textbf{filtrasyon}
$\mathbb{F} = (\mathcal{F}_n)_{n\geq0}$; $\mathcal{F}_0 \subseteq \mathcal{F}_1 \subseteq \cdots \subseteq \mathcal{F}$
artan $\sigma$-cebirleri dizisi. $\mathcal{F}_n$, zaman $n$'deki bilgiyi temsil eder.

Dizi $(X_n)_{n\geq0}$ filtrasyona \textbf{uyumlu} (adapted) dır:
her $n$ için $X_n$ $\mathcal{F}_n$-ölçülebilir.
\end{tanim}

\begin{tanim}[Doğal Filtrasyon]
$X_0, X_1, \ldots$ dizisi için \textbf{doğal filtrasyon}:
$\mathcal{F}_n = \sigma(X_0, X_1, \ldots, X_n)$ — "ilk $n+1$ gözlemden üretilen bilgi".
\end{tanim}

% =====================================================================
\section{Martingal Tanımı ve Örnekler}
\label{sec:martingal_tanim}
% =====================================================================

\begin{tanim}[Martingal]
\label{def:martingal}
$(\mathbb{F},P)$-uyumlu ve $E[|X_n|]<\infty$ olan $(X_n)$ dizisi \textbf{martingal}dır, eğer
\[
  E[X_{n+1} \mid \mathcal{F}_n] = X_n \quad \text{n.k., her } n\geq0.
\]
\begin{itemize}
  \item \textbf{Süpmartingal (supermartingale):} $E[X_{n+1}\mid\mathcal{F}_n] \leq X_n$.
  \item \textbf{Altmartingal (submartingale):} $E[X_{n+1}\mid\mathcal{F}_n] \geq X_n$.
\end{itemize}
Sezgi: martingal $\equiv$ dürüst oyun; süpmartingal $\equiv$ oyuncuya karşı oyun.
\end{tanim}

\begin{ornek}[Temel Martingal Örnekleri]
\begin{enumerate}[(1)]
  \item \textbf{Kümülatif toplam:} $X_i$ bağımsız, $E[X_i]=0$; $S_n = \sum_{i=1}^n X_i$ martingal.
        $E[S_{n+1}\mid\mathcal{F}_n] = S_n + E[X_{n+1}\mid\mathcal{F}_n] = S_n + 0 = S_n$.

  \item \textbf{Ürün martingali:} $X_i$ bağımsız, $E[X_i]=1$; $M_n = \prod_{i=1}^n X_i$ martingal.

  \item \textbf{Koşullu beklenti:} $Y \in L^1$; $M_n = E[Y\mid\mathcal{F}_n]$ martingal
        (kule özelliği: $E[M_{n+1}\mid\mathcal{F}_n] = E[E[Y\mid\mathcal{F}_{n+1}]\mid\mathcal{F}_n]
        = E[Y\mid\mathcal{F}_n] = M_n$).

  \item \textbf{Doğrusal kazanç:} $M_n = S_n^2 - n\sigma^2$ ($X_i$ bağımsız, $E[X_i]=0$, $\mathrm{Var}(X_i)=\sigma^2$) martingal.
\end{enumerate}
\end{ornek}

\begin{onerme}[Jensen — Martingalden Altmartingale]
$M_n$ martingal ve $\varphi$ konveks ise $\varphi(M_n)$ altmartingal:
$E[\varphi(M_{n+1})\mid\mathcal{F}_n] \geq \varphi(E[M_{n+1}\mid\mathcal{F}_n]) = \varphi(M_n)$.
Özel: $|M_n|$ altmartingal ($\varphi = |\cdot|$).
\end{onerme}

% =====================================================================
\section{Durma Zamanları}
\label{sec:durma_zamani}
% =====================================================================

\begin{tanim}[Durma Zamanı]
\label{def:durma_zamani}
$\tau : \Omega \to \{0,1,2,\ldots\} \cup \{\infty\}$ rasgele değişkeni
\textbf{durma zamanı} (stopping time) dır, eğer her $n$ için
$\{\tau \leq n\} \in \mathcal{F}_n$ ise (durma kararı geleceği gerektirmez).
\end{tanim}

\begin{ornek}[Durma Zamanı Örnekleri]
\begin{itemize}
  \item İlk geçiş zamanı: $\tau_a = \inf\{n \geq 0 : S_n \geq a\}$ durma zamanıdır.
  \item Sabit $\tau = n_0$: her sabit zaman durma zamanıdır.
  \item "Son $n$ için $X_n > a$": bu gelecek bilgisi gerektirir — durma zamanı \emph{değildir}.
\end{itemize}
\end{ornek}

% =====================================================================
\section{İsteğe Bağlı Durma Teoremi}
\label{sec:ibd_teoremi}
% =====================================================================

\begin{teorem}[Doob'un İsteğe Bağlı Durma Teoremi]
\label{thm:ibd}
$(M_n)$ martingal, $\tau$ durma zamanı. Aşağıdaki koşullardan biri sağlanırsa
$E[M_\tau] = E[M_0]$:
\begin{enumerate}[(i)]
  \item $\tau$ sınırlı: $\tau \leq N$ sabit $N$ için.
  \item $\tau < \infty$ n.k.\ ve $|M_n|$ düzgün integrallenebilir.
  \item $E[\tau] < \infty$ ve $|M_{n+1}-M_n| \leq C$ sabit $C$.
\end{enumerate}
\end{teorem}

\begin{proof}[Sınırlı Durum İspat]
$\tau \leq N$. $M_\tau = M_N - \sum_{k=\tau}^{N-1}(M_{k+1}-M_k) \cdot \mathbf{1}_{\tau \leq k}$...
Alternatif: dürztaması takası.
$E[M_\tau] = \sum_{n=0}^N E[M_n \mathbf{1}_{\tau=n}] = \sum_{n=0}^N E[E[M_N\mid\mathcal{F}_n]\mathbf{1}_{\tau=n}]
= E[M_N] = E[M_0]$.
($\{\tau=n\} \in \mathcal{F}_n$; martingal özelliği ve kule.)
\end{proof}

\begin{ornek}[Gamblerın İflası]
Simetrik yürüyüş: $S_n = \sum_{i=1}^n \xi_i$ ($\xi_i = \pm1$ eşit olasılıkla).
$\tau = \inf\{n: S_n = a \text{ veya } S_n = -b\}$ ($a,b > 0$).
$M_n = S_n$ martingal; IBD: $E[S_\tau] = S_0 = 0$.
$E[S_\tau] = a \cdot P(S_\tau=a) - b \cdot P(S_\tau=-b) = 0$.
$P(S_\tau=a) = b/(a+b)$, $P(S_\tau=-b) = a/(a+b)$.
Başlangıç $S_0 = 0$: $a$ liraya ulaşma olasılığı $= b/(a+b)$.
\end{ornek}

% =====================================================================
\section{Doob'un Yakınsama Teoremi}
\label{sec:doob_yaklinsama}
% =====================================================================

\begin{tanim}[Yukarı Geçiş Sayısı]
$[a,b]$ aralığı için $(X_n)$'nin $n$ adımda $[a,b]$'yi \emph{aşağıdan yukarıya}
kaç kez geçtiği sayısına $U_n(a,b)$ denir (upcrossing number).
\end{tanim}

\begin{lemma}[Doob Yukarı Geçiş Lemması]
\label{lem:doob_upcrossing}
$(X_n)$ altmartingal ve $a < b$. O zaman:
$(b-a) E[U_n(a,b)] \leq E[(X_n-a)^+] - E[(X_0-a)^+]$.
\end{lemma}

\begin{teorem}[Doob Martingal Yakınsama Teoremi]
\label{thm:doob_yaklinsama}
$(X_n)$ altmartingal, $\sup_n E[X_n^+] < \infty$. O zaman
$X_\infty = \lim_{n\to\infty} X_n$ n.k.\ var ve $E[|X_\infty|] < \infty$.
\end{teorem}

\begin{proof}
Her $a < b$ için yukarı geçiş lemması:
$(b-a)E[U_\infty(a,b)] \leq \sup_n E[(X_n-a)^+] < \infty$.
$P(U_\infty(a,b) = \infty) = 0$ her rasyonel $a < b$ için.
$P(\liminf X_n < a < b < \limsup X_n) = 0$ tüm rasyonel çiftlerde.
Sayılabilir birleşim: $P(\liminf X_n < \limsup X_n) = 0$,
yani $X_n$ her yerde limit sahibi (sonsuz olabilir).
$\sup E[X_n^+] < \infty$ $\Rightarrow$ $E[X_\infty^+] < \infty$; monoton sınır $\Rightarrow$ sonlu.
\end{proof}

\begin{onerme}[GBSK — Martingal İspatı]
$X_i$ i.i.d., $E[|X_1|] < \infty$. $S_n = \sum X_i$.
$M_n = E[S_N/N \mid \mathcal{F}_n]$ (geri-martingal argümanıyla) $\to E[X_1]$ n.k.
Bu, GBSK'nın martingal teorisi aracılığıyla zarif alternatif ispatıdır.
\end{onerme}

% =====================================================================
\section{$L^2$ Martingalleri}
\label{sec:l2_martingal}
% =====================================================================

\begin{teorem}[$L^2$ Martingal Yakınsaması]
\label{thm:l2_martingal}
$(M_n)$ martingal, $\sup_n E[M_n^2] < \infty$. O zaman
$M_\infty$ mevcut, $E[M_\infty^2] < \infty$, ve $M_n \xrightarrow{L^2} M_\infty$.
\end{teorem}

\begin{proof}
$M_n^2$ altmartingal (Jensen); $\sup E[M_n^2] < \infty$ $\Rightarrow$ Doob yakınsama teoremi.
$L^2$ yakınsama: Doob'un eşitsizliği $E[\sup_n M_n^2] \leq 4\sup E[M_n^2]$;
DCT ile $L^2$ yakınsama.
\end{proof}

\begin{ornek}[Rassal Yürüyüş — $L^2$ Martingal]
$S_n = \sum_{i=1}^n \xi_i$ ($E[\xi_i]=0$, $E[\xi_i^2]=\sigma^2$).
$M_n = S_n/\sqrt{n}$: $E[M_n^2] = \sigma^2$ sınırlı.
Ama $M_n/\sqrt{\ln n} \to 0$ n.k.\ (YLK'nın sonucu),
$M_n$ ise $\Normal(0,\sigma^2)$'e dağılımla yakınsar (MLT).
$L^2$ sınırlılık n.k.\ yakınsamayı garanti etmez (sınır değişkenlik gösterebilir).
\end{ornek}

% ---- Öz-Yansıtma ----
\begin{ozyansima}
\begin{itemize}[leftmargin=1.8em]
  \item "Martingal = dürüst oyun" metaforunu iki farklı örnekle somutlaştırın.
  \item IBD teoremi neden durma zamanının sınırlı olmasını (veya başka bir koşul) gerektiriyor?
  \item GBSK'nın martingal ispatı neden daha "zarif"tir?
\end{itemize}
\end{ozyansima}

% ---- Köprü Notu ----
\begin{kopru}
\textbf{Önceki:} Bölüm~\ref{sec:kosullu_beklenti} koşullu beklenti; Böl.~\ref{chp:buyuk_sayilar} BSK.\\
\textbf{Sonraki (opsiyonel):} Bölüm~\ref{chp:stokastik} — Stokastik Süreçler.\\
\textbf{Uygulamalar:} Finansta martingal fiyatlama; sıralı analizde isteğe bağlı durma.
\defterbaglantisi{bolum09\_martingaller.ipynb}{Martingal simülasyonu, gamblerın iflası}
\end{kopru}

% =====================================================================
%  ALIŞTIRMALAR
% =====================================================================

\cozumlualistirmalar

\begin{alistirma}[Al.9.1]{A}{Martingal Doğrulama}
$X_i$ bağımsız, $E[X_i]=1$, $M_n = \prod_{i=1}^n X_i$.
$E[M_{n+1}\mid\mathcal{F}_n] = M_n$ olduğunu gösterin.
\end{alistirma}
\alcozum{%
$E[M_{n+1}\mid\mathcal{F}_n] = E[M_n X_{n+1}\mid\mathcal{F}_n] = M_n E[X_{n+1}\mid\mathcal{F}_n]
= M_n E[X_{n+1}] = M_n \cdot 1 = M_n$.
(Bağımsızlık: $X_{n+1}$ $\mathcal{F}_n$'den bağımsız.)}

\begin{alistirma}[Al.9.2]{B}{Gamblerın İflası — Varyans Martingali}
$S_n$ simetrik rassal yürüyüş ($\xi_i = \pm1$, $\sigma^2=1$). $M_n = S_n^2 - n$.
(a) $M_n$'nin martingal olduğunu gösterin.
(b) $\tau = \inf\{n: S_n = a\}$ ($a > 0$) için $E[\tau] = a^2$ olduğunu IBD ile bulun.
\end{alistirma}
\alcozum{%
(a) $E[M_{n+1}\mid\mathcal{F}_n] = E[(S_n+\xi_{n+1})^2 - (n+1)\mid\mathcal{F}_n]
= S_n^2 + 2S_n E[\xi_{n+1}] + E[\xi_{n+1}^2] - n - 1 = S_n^2 - n = M_n$. $\checkmark$\\
(b) IBD ($\tau$ sınırsız ama Koşul (iii) çalışır — artışlar $\leq 1+1+1$):
$E[M_\tau] = E[S_\tau^2 - \tau] = E[M_0] = 0$.
$S_\tau = a$ (emici bariyer): $a^2 - E[\tau] = 0$. $E[\tau] = a^2$.}

% ---

\alistirmalar

\begin{alistirma}[Al.9.3]{A}{Altmartingal}
$(X_n)$ martingal. $\varphi(x) = x^2$ için $\varphi(X_n)$'nin altmartingal
olduğunu Jensen'i kullanarak gösterin.
\end{alistirma}
\alcozum{%
$E[X_{n+1}^2\mid\mathcal{F}_n] \geq (E[X_{n+1}\mid\mathcal{F}_n])^2 = X_n^2$. (Jensen, $\varphi$ konveks.)}

\begin{alistirma}[Al.9.4]{B}{Doob Yakınsama Uygulaması}
$X_i \geq 0$ bağımsız, $E[X_i] = 1$. $M_n = \prod_{i=1}^n X_i$.
$\sup_n E[M_n] = 1 < \infty$ ise $M_n$ n.k.\ yakınsar.
$M_\infty = 0$ n.k.\ olduğunu BSK ile gösterin.
\end{alistirma}
\alcozum{%
$\ln M_n = \sum \ln X_i$. BSK: $(\ln M_n)/n \to E[\ln X_1]$.
Jensen: $E[\ln X_1] \leq \ln E[X_1] = 0$; genellikle $< 0$.
$\ln M_n / n \to c < 0$ $\Rightarrow$ $M_n = e^{n \cdot c+o(n)} \to 0$ n.k.}

% =====================================================================
\endinput
